% ===================================================================
\section{The Reduced Axiom Basis}
\label{ch:basis}
% ===================================================================

\subsection{From 11 Premises to 7}

The original published FRFP theory enumerates 11 named premises: ETS, RB, AE, MD,
HEG, HEC, CP, IL, AR, NTER, CSC. The Lean formalization establishes that four of
these (ETS, RB, AE, MD) are theorems derivable directly from the kernel
definitions. The minimal independent basis therefore has exactly \textbf{7
primitive assumptions}.

\begin{center}
\begin{tabular}{llll}
\toprule
Premise & Original status & Lean status & Location \\
\midrule
ETS  & Axiom & \textbf{Theorem} (kernel typing) & Thm 2.1 \\
RB   & Axiom & \textbf{Theorem} (kernel enumeration) & Thm 2.2 \\
AE   & Axiom & \textbf{Theorem} (AI-executable typing) & Thm 2.3 \\
MD   & Axiom & \textbf{Theorem} (AI/human partition) & Thm 2.4 \\
HEG  & Axiom & \textbf{Axiom} (retained) & \S\ref{sec:core-irreducible} \\
HEC  & Axiom & \textbf{Axiom} (retained) & \S\ref{sec:core-irreducible} \\
CP   & Axiom & \textbf{Axiom} (retained) & \S\ref{sec:core-irreducible} \\
IL   & Axiom & \textbf{Axiom} (retained) & \S\ref{sec:interaction} \\
AR   & Axiom & \textbf{Axiom} (retained) & \S\ref{sec:interaction} \\
NTER & Axiom & \textbf{Axiom} (retained) & \S\ref{sec:interaction} \\
CSC  & Axiom & \textbf{Axiom} (retained) & \S\ref{sec:interaction} \\
\bottomrule
\end{tabular}
\end{center}

% -------------------------------------------------------------------
\subsection{Derivable Theorems (No Longer Axioms)}
% -------------------------------------------------------------------

\subsubsection*{ETS --- Explicit--Tacit Separation}

\begin{theorem}[ETS is a theorem]
There is no primitive from $\Tobj$ to $\Eobj$:
\[
  \forall p : \mathrm{Primitive},\
  \neg(p.\mathrm{source} = \Tobj \wedge p.\mathrm{target} = \Eobj)
\]
\leanref{Frfp.Minimal.ReducedBasis.ETS\_derivable} --- proved by delegating to
\texttt{no\_morphism\_T0\_to\_E0}.
\end{theorem}

\begin{note}
ETS holds because the primitive enumeration simply contains no $\Tobj \to \Eobj$
morphism. It is a fact about the closed enumeration, not an independent assumption.
\end{note}

\subsubsection*{RB --- Boundary Uniqueness}

\begin{theorem}[RB uniqueness is a theorem]
$\RB$ is the unique primitive from $\Eobj$ to $\Tobj$:
\[
  \forall p : \mathrm{Primitive},\
  (p.\mathrm{source} = \Eobj \wedge p.\mathrm{target} = \Tobj) \Rightarrow p = \RB
\]
\leanref{Frfp.Minimal.ReducedBasis.RB\_derivable}
\end{theorem}

\subsubsection*{AE --- AI-Explicit Restriction}

\begin{theorem}[AE is a theorem]
Every AI-executable primitive stays within explicit space:
\[
  \forall p,\
  \textit{is\_AI\_executable}(p)
  \Rightarrow (p.\mathrm{source} \in \{\varnothing, \Eobj\})
              \wedge p.\mathrm{target} = \Eobj
\]
\leanref{Frfp.Minimal.ReducedBasis.AE\_derivable}
\end{theorem}

\subsubsection*{MD --- Mode Discipline}

\begin{theorem}[MD is a theorem]
AI-executable primitives always target $\Eobj$; human-only primitives always
target $\Tobj$:
\[
  (\forall p,\ \textit{is\_AI\_executable}(p) \Rightarrow p.\mathrm{target} = \Eobj)
  \wedge
  (\forall p,\ \textit{is\_human\_only}(p) \Rightarrow p.\mathrm{target} = \Tobj)
\]
\leanref{Frfp.Minimal.ReducedBasis.MD\_derivable}
\end{theorem}

% -------------------------------------------------------------------
\subsection{Core Irreducible Premises}
\label{sec:core-irreducible}
% -------------------------------------------------------------------

The three core premises below are retained as genuine axioms. They capture
domain commitments that cannot be derived from the kernel structure alone.

\subsubsection*{HEG --- Human-Exclusive Grounding}

\begin{axiom_env}[Human-Exclusive Grounding (HEG)]
Tacit evaluation and human final decision --- the primitives $\TE$ and $\HFD$ ---
are exclusively executable by humans. No AI system can perform tacit grounding.

\emph{Formal skeleton}:
$\mathrm{HEG}_{\mathrm{skel}}:\ \TE.\mathrm{target} = \Tobj \wedge
\HFD.\mathrm{target} = \Tobj$

\leanref{Frfp.Minimal.ReducedBasis.HEG\_skeleton\_derivable}

\emph{Policy residue}: The normative claim --- that only humans may execute
$\TE$ and $\HFD$ --- is a governance premise.
\end{axiom_env}

\subsubsection*{HEC --- Human-Exclusive Closure}

\begin{axiom_env}[Human-Exclusive Closure (HEC)]
The final-decision primitive $\HFD$ is an endomorphism of tacit space that
only humans may apply.

\emph{Formal skeleton}:
$\mathrm{HEC}_{\mathrm{skel}}:\ \HFD.\mathrm{source} = \Tobj \wedge
\HFD.\mathrm{target} = \Tobj$

\leanref{Frfp.Minimal.ReducedBasis.HEC\_skeleton\_derivable}
\end{axiom_env}

\subsubsection*{CP --- Compositional Pipelines}

\begin{axiom_env}[Compositional Pipelines (CP)]
Pipelines compose in a way that preserves the source and target boundary:
\[
  (\Pi_2 \circ \Pi_1).\mathrm{source} = \Pi_1.\mathrm{source},\quad
  (\Pi_2 \circ \Pi_1).\mathrm{target} = \Pi_2.\mathrm{target}
\]
\leanref{Frfp.Minimal.ReducedBasis.CP\_skeleton\_derivable}
\end{axiom_env}

% -------------------------------------------------------------------
\subsection{Interaction Irreducible Premises}
\label{sec:interaction}
% -------------------------------------------------------------------

\subsubsection*{IL --- Intent Locking}

\begin{axiom_env}[Intent Locking (IL)]
All AI steps in a protocol trace share a common intent tag. An AI step cannot
unilaterally change the declared intent of the interaction:
\[
  \forall a, b \in \mathcal{P},\
  a.\mathrm{isAI} = \top \wedge b.\mathrm{isAI} = \top
  \Rightarrow a.\mathrm{intentTag} = b.\mathrm{intentTag}
\]
\emph{Source anchor}: Temporal safety / invariant-over-traces semantics.
\textbf{[Pnueli 1977]}
\end{axiom_env}

\subsubsection*{AR --- Ambiguity Resolution}

\begin{axiom_env}[Ambiguity Resolution (AR)]
Every ambiguous step in a protocol trace must be clarified before the
interaction proceeds:
\[
  \forall a \in \mathcal{P},\
  a.\mathrm{ambiguous} = \top \Rightarrow a.\mathrm{clarified} = \top
\]
\emph{Source anchor}: Liveness over traces; safety/liveness decomposition.
\textbf{[Alpern--Schneider 1985]}
\end{axiom_env}

\subsubsection*{NTER --- No Tacit Emulation/Reconstruction}

\begin{axiom_env}[No Tacit Emulation or Reconstruction (NTER)]
No AI step attempts to construct a morphism from $\Tobj$ to $\Eobj$, i.e.,
no AI step attempts to emulate or substitute for tacit human judgment.

\emph{Formal skeleton}:
$\forall p,\ p.\mathrm{source} = \Tobj \Rightarrow p.\mathrm{target} \neq \Eobj$

\leanref{Frfp.Minimal.ReducedBasis.NTER\_skeleton\_derivable}
\end{axiom_env}

\subsubsection*{CSC --- Correctness Stability Constraint}

\begin{axiom_env}[Correctness Stability Constraint (CSC)]
Once a step in a protocol trace is correctness-locked, all subsequent steps are
correctness-locked. Correctness locking is monotone:
\[
  (\exists a \in \mathcal{P},\ a.\mathrm{correctnessLocked} = \top)
  \Rightarrow
  \forall b \in \mathcal{P},\ b.\mathrm{correctnessLocked} = \top
\]
\emph{Source anchor}: Temporal safety (invariant preservation).
\textbf{[Pnueli 1977]}
\end{axiom_env}

% -------------------------------------------------------------------
\subsection{Source Anchoring of External Assumptions}
% -------------------------------------------------------------------

Modules beyond the kernel require additional axioms from established mathematics.

\begin{center}
\begin{tabular}{lll}
\toprule
Family & Representative axioms & Primary source \\
\midrule
IEEE 754 float arithmetic & Associativity, commutativity, order & [IEEE 754-2019] \\
Probability / measure & Countable additivity, normalization & [Billingsley 1995] \\
Stopping times & Measurability, Markov property & [Billingsley 1995] \S36 \\
Stochastic tail bounds & Markov inequality, geometric decay & [Durrett 2019] \\
Term rewriting / ARS & Termination, confluence & [Baader--Nipkow 1998] \\
Newman's Lemma & Confluence from termination & [Newman 1942] \\
Category / groupoid & Associativity, identity, inverses & [Mac Lane 1971] \\
Real analysis & Geometric limit-to-zero & [Rudin 1976] \\
Temporal safety & Safety properties over traces & [Pnueli 1977] \\
Safety vs liveness & Decomposition of liveness & [Alpern--Schneider 1985] \\
\bottomrule
\end{tabular}
\end{center}

\leanref{Frfp.Minimal.ReducedBasis.sourceBackedFamilies}

% -------------------------------------------------------------------
\subsection{Decomposition: Skeleton vs Policy Residue}
% -------------------------------------------------------------------

Each of the 7 retained axioms decomposes into two parts:
\begin{itemize}
  \item \textbf{Skeleton}: the mathematically structural part, derivable or
        source-anchored to published mathematics.
  \item \textbf{Policy residue}: the irreducible normative/governance commitment.
\end{itemize}

\begin{center}
\begin{tabular}{p{1cm}p{5.5cm}p{5.5cm}}
\toprule
Axiom & Skeleton & Policy residue \\
\midrule
HEG  & $\TE, \HFD$ target $\Tobj$ (Lean-proved) & Only humans may execute these \\
HEC  & $\HFD$ is a $\Tobj$-endomorphism (Lean-proved) & No AI substitution for closure \\
CP   & Composition preserves source/target (Lean-proved) & Pipelines are the only valid composition \\
IL   & Trace invariant (Pnueli 1977) & Intent is a shared, locked resource \\
AR   & Liveness (Alpern--Schneider 1985) & Ambiguity must be resolved \\
NTER & Structural = ETS (Lean-proved) & AI prohibited from tacit emulation attempts \\
CSC  & Safety/monotonicity (Pnueli 1977) & Correctness lock is irreversible \\
\bottomrule
\end{tabular}
\end{center}

The policy residue of all seven axioms is where the human-authority commitments
of FRFP live. These cannot be formalized as mathematical theorems --- they are
governance premises.

This is also the point at which human validation applies in the FRFP workflow.
Humans are not expected to re-verify Lean proof steps; they validate whether these
residues are acceptable commitments for the intended deployment context.

\leanref{Frfp.Minimal.ReducedBasis.DecompositionStatus}
